% --- Archon physics formalization source begin ---
% archon:physics
% archon:covers PhyXMiniProblems/problem_phyx_mini_0542.lean
% archon:source-report reports/phyx_mini/problem_phyx_mini_0542.source.json

\chapter{Physics problem phyx\_mini\_0542}
\label{ch:PhyXMiniProblems_problem_phyx_mini_0542}

\paragraph{Problem source.}
\#\# Physical scenario

A beam of spin-\textbackslash{}(1/2\textbackslash{}) particles is sent through a series of three Stern-Gerlach analyzers, as shown in figure. The second Stern-Gerlach analyzer is aligned along the \textbackslash{}(\textbackslash{}hat\{n\}\textbackslash{}) direction, which makes an angle \textbackslash{}(\textbackslash{}theta\textbackslash{}) in the \textbackslash{}(x\textbackslash{})-\textbackslash{}(z\textbackslash{}) plane with respect to the \textbackslash{}(z\textbackslash{})-axis.

\#\# Question

Find the probability that particles transmitted through the first Stern-Gerlach analyzer are measured to have spin down at the third Stern-Gerlach analyzer?

\#\# Answer choices

- A: A: \textbackslash{}(\textbackslash{}frac\{1\}\{4\} \textbackslash{}cot\textasciicircum{}2 \textbackslash{}theta\textbackslash{})

- B: B: \textbackslash{}(\textbackslash{}frac\{1\}\{4\} \textbackslash{}tan\textasciicircum{}2 \textbackslash{}theta\textbackslash{})

- C: C: \textbackslash{}(\textbackslash{}frac\{1\}\{4\} \textbackslash{}cos\textasciicircum{}2 \textbackslash{}theta\textbackslash{})

- D: D: \textbackslash{}(\textbackslash{}frac\{1\}\{4\} \textbackslash{}sin\textasciicircum{}2 \textbackslash{}theta\textbackslash{})

\#\# Figure caption (auxiliary; use the image as primary evidence)

The image depicts a flowchart or diagram with a sequence of interconnected components. Here's a breakdown of the elements:

1. **Left Rectangle with Lines**: On the far left, there is a rectangle with horizontal lines extending towards the right. This might represent an input or data source.

2. **First Component (Z Block)**: 
   - A rectangular block labeled with a large "Z" inside.
   - The block has arrows pointing upwards and downwards on its right side, indicating possible flow directions or operations.

3. **Second Component (h Block)**:
   - A rectangular block labeled with a large "h" inside.
   - Similar to the previous block, it has arrows pointing upwards and downwards on its right side.

4. **Third Component (Z Block)**:
   - Another rectangular block labeled with a large "Z."
   - It also features arrows pointing upwards and downwards on its right side.

5. **Two Output Paths**:
   - After the third block, the diagram splits into two separate paths.
   - Each path leads to a rectangular box with a question mark. These boxes likely represent unknown or variable outputs to be determined.

The combination of these elements suggests a process or system where data flows from left to right through a series of operations indicated by the blocks, with the outcomes branching into two possibilities.

\#\# Dataset metadata

Quantum Phenomena; Physical Model Grounding Reasoning, Multi-Formula Reasoning

\paragraph{Recorded answer/context.}
D: D: \textbackslash{}(\textbackslash{}frac\{1\}\{4\} \textbackslash{}sin\textasciicircum{}2 \textbackslash{}theta\textbackslash{})

\paragraph{Figure/image path.}
/root/proposal\_for\_physic/science-mango/phyx\_mini\_run/phyx\_data/test\_image/542.png

\paragraph{Formalization target.}
create a compiling Lean file with sorry bodies at `PhyXMiniProblems/problem\_phyx\_mini\_0542.lean`. The Lean declarations must preserve the physical quantities, dimensions or dimensional roles, figure labels, governing-law hypotheses, and final relation expressed by this problem.
Use Mathlib/Physlib names found through LeanExplore where available. If a physics API is missing, introduce faithful local abstractions rather than scalar placeholder aliases.

\begin{theorem}[Physics formalization target]
\label{thm:physics:phyx_mini_0542:target}
\lean{PhyXMiniProblems.ProblemPhyXMini0542.final_spin_down_probability_and_recorded_answer_D}
\uses{def:physics:phyx-mini-0542:phyxminiproblems-problemphyxmini0542-sterngerlachexperiment, def:physics:phyx-mini-0542:phyxminiproblems-problemphyxmini0542-matchesspinhalfscenario, def:physics:phyx-mini-0542:phyxminiproblems-problemphyxmini0542-matchessuppliedthreeanalyzerfigure, def:physics:phyx-mini-0542:phyxminiproblems-problemphyxmini0542-secondanalyzergeometry, def:physics:phyx-mini-0542:phyxminiproblems-problemphyxmini0542-spinhalfbornrule, def:physics:phyx-mini-0542:phyxminiproblems-problemphyxmini0542-satisfiessequentialfilteringlaw, def:physics:phyx-mini-0542:phyxminiproblems-problemphyxmini0542-iscorrectanswer}
The assigned autoformalize agent should translate this physics problem into Lean declarations in the covered file, with theorem and lemma proof bodies written as `by sorry`.
\end{theorem}
\begin{proof}
This is an autoformalization task, not a proof task. Produce statements that can later be proved by the physics prover without weakening the physical model.
\end{proof}

% --- Archon declaration topology begin ---
\section{Declaration topology}
\begin{definition}[Spin Outcome]
  \label{def:physics:phyx-mini-0542:phyxminiproblems-problemphyxmini0542-spinoutcome}
  \lean{PhyXMiniProblems.ProblemPhyXMini0542.SpinOutcome}
  The two output channels of a spin-one-half Stern–Gerlach analyzer
\end{definition}
\begin{proof}
  This is the declared model definition of Spin Outcome.
\end{proof}
\begin{definition}[opposite]
  \label{def:physics:phyx-mini-0542:phyxminiproblems-problemphyxmini0542-spinoutcome-opposite}
  \lean{PhyXMiniProblems.ProblemPhyXMini0542.SpinOutcome.opposite}
  \uses{def:physics:phyx-mini-0542:phyxminiproblems-problemphyxmini0542-spinoutcome}
  Reversing the spin outcome along a fixed oriented analyzer axis
\end{definition}
\begin{proof}
  This is the declared model definition of opposite.
\end{proof}
\begin{definition}[Particle Spin]
  \label{def:physics:phyx-mini-0542:phyxminiproblems-problemphyxmini0542-particlespin}
  \lean{PhyXMiniProblems.ProblemPhyXMini0542.ParticleSpin}
  The intrinsic spin quantum number specified for the incident particles
\end{definition}
\begin{proof}
  This is the declared model definition of Particle Spin.
\end{proof}
\begin{definition}[Spin Axis]
  \label{def:physics:phyx-mini-0542:phyxminiproblems-problemphyxmini0542-spinaxis}
  \lean{PhyXMiniProblems.ProblemPhyXMini0542.SpinAxis}
  An oriented analyzer axis. Its direction points toward the channel called up, so the unit-vector orientation retains the physical distinction between spin up and spin down
\end{definition}
\begin{proof}
  This is the declared model definition of Spin Axis.
\end{proof}
\begin{definition}[dimensionless Pauli Observable]
  \label{def:physics:phyx-mini-0542:phyxminiproblems-problemphyxmini0542-spinaxis-dimensionlesspauliobservable}
  \lean{PhyXMiniProblems.ProblemPhyXMini0542.SpinAxis.dimensionlessPauliObservable}
  \uses{def:physics:phyx-mini-0542:phyxminiproblems-problemphyxmini0542-spinaxis}
  The dimensionless spin observable n · σ associated with an oriented analyzer axis n. The three matrices are PhysLean's Pauli matrices; multiplying this observable by ℏ / 2 would give the corresponding physical spin component
\end{definition}
\begin{proof}
  This is the declared model definition of dimensionless Pauli Observable.
\end{proof}
\begin{definition}[Analyzer Axis Label]
  \label{def:physics:phyx-mini-0542:phyxminiproblems-problemphyxmini0542-analyzeraxislabel}
  \lean{PhyXMiniProblems.ProblemPhyXMini0542.AnalyzerAxisLabel}
  Axis symbols printed inside the three analyzer boxes in the figure
\end{definition}
\begin{proof}
  This is the declared model definition of Analyzer Axis Label.
\end{proof}
\begin{definition}[Stern Gerlach Analyzer]
  \label{def:physics:phyx-mini-0542:phyxminiproblems-problemphyxmini0542-sterngerlachanalyzer}
  \lean{PhyXMiniProblems.ProblemPhyXMini0542.SternGerlachAnalyzer}
  \uses{def:physics:phyx-mini-0542:phyxminiproblems-problemphyxmini0542-spinaxis, def:physics:phyx-mini-0542:phyxminiproblems-problemphyxmini0542-analyzeraxislabel}
  A Stern–Gerlach analyzer together with its displayed axis label
\end{definition}
\begin{proof}
  This is the declared model definition of Stern Gerlach Analyzer.
\end{proof}
\begin{definition}[Stern Gerlach Experiment]
  \label{def:physics:phyx-mini-0542:phyxminiproblems-problemphyxmini0542-sterngerlachexperiment}
  \lean{PhyXMiniProblems.ProblemPhyXMini0542.SternGerlachExperiment}
  \uses{def:physics:phyx-mini-0542:phyxminiproblems-problemphyxmini0542-spinoutcome, def:physics:phyx-mini-0542:phyxminiproblems-problemphyxmini0542-particlespin, def:physics:phyx-mini-0542:phyxminiproblems-problemphyxmini0542-sterngerlachanalyzer}
  The three-analyzer experiment and its requested dimensionless probability. probabilityConditionedOnFirstTransmission is an independent observable: it is the probability, among particles emerging through the selected output of the first analyzer, of following the selected output of the second analyzer and then emerging through the requested output of the third analyzer
\end{definition}
\begin{proof}
  This is the declared model definition of Stern Gerlach Experiment.
\end{proof}
\begin{definition}[Matches Spin Half Scenario]
  \label{def:physics:phyx-mini-0542:phyxminiproblems-problemphyxmini0542-matchesspinhalfscenario}
  \lean{PhyXMiniProblems.ProblemPhyXMini0542.MatchesSpinHalfScenario}
  \uses{def:physics:phyx-mini-0542:phyxminiproblems-problemphyxmini0542-sterngerlachexperiment}
  Scenario facts stated in the prose, including the physical probability range
\end{definition}
\begin{proof}
  This is the declared model definition of Matches Spin Half Scenario.
\end{proof}
\begin{definition}[Matches Supplied Three Analyzer Figure]
  \label{def:physics:phyx-mini-0542:phyxminiproblems-problemphyxmini0542-matchessuppliedthreeanalyzerfigure}
  \lean{PhyXMiniProblems.ProblemPhyXMini0542.MatchesSuppliedThreeAnalyzerFigure}
  \uses{def:physics:phyx-mini-0542:phyxminiproblems-problemphyxmini0542-sterngerlachexperiment}
  Analyzer labels, coordinate directions, and connected output ports read from the supplied image. Coordinates 0, 1, and 2 represent x, y, and z. The upper channels of the first two analyzers continue to the next box, while the question asks about the lower (down) channel of the last box
\end{definition}
\begin{proof}
  This is the declared model definition of Matches Supplied Three Analyzer Figure.
\end{proof}
\begin{definition}[Second Analyzer Geometry]
  \label{def:physics:phyx-mini-0542:phyxminiproblems-problemphyxmini0542-secondanalyzergeometry}
  \lean{PhyXMiniProblems.ProblemPhyXMini0542.SecondAnalyzerGeometry}
  \uses{def:physics:phyx-mini-0542:phyxminiproblems-problemphyxmini0542-sterngerlachexperiment}
  The second analyzer's unit axis lies in the x-z plane and makes the dimensionless angle theta, read in radians, with the first positive z axis. The usual range of the undirected Euclidean angle is recorded explicitly
\end{definition}
\begin{proof}
  This is the declared model definition of Second Analyzer Geometry.
\end{proof}
\begin{definition}[Spin Half Born Rule]
  \label{def:physics:phyx-mini-0542:phyxminiproblems-problemphyxmini0542-spinhalfbornrule}
  \lean{PhyXMiniProblems.ProblemPhyXMini0542.SpinHalfBornRule}
  \uses{def:physics:phyx-mini-0542:phyxminiproblems-problemphyxmini0542-spinoutcome, def:physics:phyx-mini-0542:phyxminiproblems-problemphyxmini0542-spinaxis}
  Born transition probabilities for spin one-half measurements. For analyzer axes separated by an undirected angle alpha, retaining the same spin label has probability cos²(alpha / 2) and obtaining the opposite label has probability sin²(alpha / 2). The law is general in both axes and does not mention this problem's requested final probability
\end{definition}
\begin{proof}
  This is the declared model definition of Spin Half Born Rule.
\end{proof}
\begin{definition}[Satisfies Sequential Filtering Law]
  \label{def:physics:phyx-mini-0542:phyxminiproblems-problemphyxmini0542-satisfiessequentialfilteringlaw}
  \lean{PhyXMiniProblems.ProblemPhyXMini0542.SatisfiesSequentialFilteringLaw}
  \uses{def:physics:phyx-mini-0542:phyxminiproblems-problemphyxmini0542-sterngerlachexperiment, def:physics:phyx-mini-0542:phyxminiproblems-problemphyxmini0542-spinhalfbornrule}
  The chain rule for two successive ideal filters after conditioning on passage through the first analyzer. In particular, no probability for entering the first selected port is multiplied into this conditional probability
\end{definition}
\begin{proof}
  This is the declared model definition of Satisfies Sequential Filtering Law.
\end{proof}
\begin{definition}[Answer Choice]
  \label{def:physics:phyx-mini-0542:phyxminiproblems-problemphyxmini0542-answerchoice}
  \lean{PhyXMiniProblems.ProblemPhyXMini0542.AnswerChoice}
  Labels of the four answer choices printed with the problem
\end{definition}
\begin{proof}
  This is the declared model definition of Answer Choice.
\end{proof}
\begin{definition}[displayed Probability]
  \label{def:physics:phyx-mini-0542:phyxminiproblems-problemphyxmini0542-displayedprobability}
  \lean{PhyXMiniProblems.ProblemPhyXMini0542.displayedProbability}
  \uses{def:physics:phyx-mini-0542:phyxminiproblems-problemphyxmini0542-answerchoice}
  The dimensionless probability expression printed beside each answer label
\end{definition}
\begin{proof}
  This is the declared model definition of displayed Probability.
\end{proof}
\begin{definition}[Is Correct Answer]
  \label{def:physics:phyx-mini-0542:phyxminiproblems-problemphyxmini0542-iscorrectanswer}
  \lean{PhyXMiniProblems.ProblemPhyXMini0542.IsCorrectAnswer}
  \uses{def:physics:phyx-mini-0542:phyxminiproblems-problemphyxmini0542-sterngerlachexperiment, def:physics:phyx-mini-0542:phyxminiproblems-problemphyxmini0542-answerchoice, def:physics:phyx-mini-0542:phyxminiproblems-problemphyxmini0542-displayedprobability}
  A displayed choice agrees with the experiment's requested probability
\end{definition}
\begin{proof}
  This is the declared model definition of Is Correct Answer.
\end{proof}
% --- Archon declaration topology end ---
% --- Archon physics formalization source end ---
